% Original: PDF strana 24, donji slajd.
\slajdid{02-s048}
\begin{frame}[label=02-s048]{NI realizacija: tri različite vrste čipova}
\setlength{\parskip}{0pt}
Da bismo koristili kola iste vrste, funkciju najčešće algebarski preoblikujemo za realizaciju NI ili NILI kolima. Za NI realizaciju po pravilu je najlakše poći od zbira proizvoda:
% element: 02-s048:e01
\[F=\bar D\bar B+D\bar C\bar A
=\overline{\overline{\bar D\bar B+D\bar C\bar A}}
=\overline{\overline{\bar D\bar B}\cdot\overline{D\bar C\bar A}}\]
\par\vspace{2mm}
% element: 02-s048:e02
Potrebni su:
\begin{itemize}
\item jedan čip sa najmanje \textbf{četiri invertora};
\item jedan čip sa najmanje \textbf{dva dvoulazna NI kola};
\item jedan čip sa najmanje \textbf{jednim troulaznim NI kolom}.
\end{itemize}
Ukupno su to \textbf{tri čipa različitih vrsta}.
\par\vspace{3mm}
% element: 02-s048:e04
\Rezultat{NI transformacija \alert{\textbf{ne dokazuje minimalnost}} prema klasičnoj minimizaciji.}
\end{frame}
